\documentclass[a4]{article}

% PAQUETES %%%%%%%%%%%%%%%%%%%%%%%%%%%%%%%%%%%%%%%%%%%%%%%%%%%%%%%%%%%%%%%%%%%%%
\usepackage[utf8]{inputenc}
\usepackage{framed} % Para crear cuadros alrededor del texto
\usepackage{makecell}
\usepackage[english]{babel}
%
\usepackage[fixlanguage]{babelbib}
    \bibliographystyle{IEEEtran}

\usepackage{url}

%!TEX encoding = UTF-8 Unicode
% !TEX spellcheck = es

\setlength{\textwidth}{6.5 in}
\setlength{\textheight}{8.5 in}
\setlength{\headsep}{0.5 in}

\usepackage{multirow, array} % para las tablas
\usepackage{float} % para usar [H]
\usepackage{charter}

\usepackage[breaklinks=true]{hyperref}

%
% ADD PACKAGES here:
%

\usepackage[utf8]{inputenc}
\usepackage{amsmath,longtable,siunitx, float, graphicx, dcolumn, bm, fancyhdr, authblk, subcaption, caption, hyperref, multicol, setspace}

\usepackage{bbding}
\usepackage{amssymb}
\usepackage[rightcaption]{sidecap}
%\usepackage{textcomp}

\usepackage[hmarginratio=1:1, top=32mm, columnsep=20pt, headsep=\dimexpr20mm, margin=0.7in, top=1in ]{geometry}

\usepackage{booktabs}
\usepackage{multirow}

\usepackage{algorithm}
\usepackage{enumitem}
\usepackage{tabularx}
\usepackage[noend]{algpseudocode}
\usepackage{xcolor}

\usepackage{tikz}
\usetikzlibrary{shapes.geometric, arrows.meta, positioning}

% ----------- global styles -----------
\tikzset{
  font=\scriptsize,                      % tiny text everywhere
  every node/.style       = {transform shape},   % respect scale
  startstop/.style        = {ellipse, draw, fill=blue!10,
                             inner sep=1pt, minimum width=13mm},
  process/.style          = {rectangle, draw, fill=orange!15,
                             text width=26mm, inner sep=2pt, align=center},
  decision/.style         = {diamond, draw, fill=green!15, aspect=2.3,
                             text width=20mm, inner sep=1pt, align=center},
  arrow/.style            = {->, >=Stealth, shorten >=1pt, shorten <=1pt}
}

\newcommand{\TODO}{\textbf{\textcolor{red}{TODO}}}

\renewcommand{\algorithmicrequire}{\textbf{Input:}}
\renewcommand{\algorithmicensure}{\textbf{Output:}}

\lhead{\includegraphics[scale=0.4]{CETC.png}}
\rhead{\includegraphics[scale=0.025]{Cornell-University-Logo.png}}

\newtheorem{theo}{Theorem}
\newenvironment{theorem}
  {\begin{mdframed}\begin{theo}}
  {\end{theo}\end{mdframed}}
  
\newtheorem{prop}{Proposición}
\newenvironment{proposition}
  {\begin{mdframed}\begin{prop}}
  {\end{prop}\end{mdframed}}

\newtheorem{lem}{Lema}
\newenvironment{lema}
  {\begin{mdframed}\begin{lem}}
  {\end{lem}\end{mdframed}}
  
\newtheorem{defi}{Definition}
\newenvironment{definition}
  {\begin{mdframed}\begin{def}}
  {\end{def}\end{mdframed}}
  
\newtheorem{proof}{Proof}
\newtheorem{remark}{Remark}

\usepackage[numbers]{natbib}
\renewcommand{\thesection}{\Roman{section}}
%\setlength{\parskip}{\baselineskip}%

\begin{document}

\thispagestyle{fancy}

\begin{center}
    \textbf{\large Orbit Odyssey - Computer Vision} \\
    {Quotation for phase 1}
    
    \vspace{0.7em}

    David Alejandro Fuquen \\

    {\small \today}

\end{center}

\section*{Quotation for Phase 1 of the Orbit Odyssey - Computer Vision Autonomous Challenge}


\subsection*{Completed Work}

\begin{enumerate}[leftmargin=*]
    \item \textbf{Playing Field Definition} (3.5 hrs)
    \begin{itemize}
        \item Field dimensions, boundaries, and markings
        \item Robot starting zone placement rules
        \item Object quantity, placement templates, and clearance constraints
    \end{itemize}

    \item \textbf{Scoring System and Scorecards} (3.5 hrs)
    \begin{itemize}
        \item Complete point system (field performance, mAP bonus, model-size multiplier)
        \item Penalty definitions and edge-case rulings
        \item Tiebreaker hierarchy
        \item Printable scorecard template for referees
        \item Integration of cryptographic receipt verification into scoring workflow (see Phase 2 section for details)
    \end{itemize}

    \item \textbf{Student Guide} (2 hrs)
    \begin{itemize}
        \item Step-by-step instructions for using the training GUI
        \item Explanation of hyperparameter choices and their trade-offs
        \item How to submit the trained model and receipt for deployment
    \end{itemize}

    \item \textbf{Moderator Guide} (2 hrs)
    \begin{itemize}
        \item Field setup checklist (per-match randomization)
        \item Receipt verification procedure (decryption and validation) (see Phase 2 section for details)
    \end{itemize}

    \item \textbf{Anti-Cheating Documentation} (2 hr)
    \begin{itemize}
        \item How the signed receipt system works (for judges)
        \item What constitutes a valid vs.\ tampered receipt
        \item Disqualification criteria
    \end{itemize}

        \item \textbf{Final revision and corrections} (2 hr)
    \begin{itemize}
        \item Synthesis of work into the final document, shared via google drive
    \end{itemize}
\end{enumerate}

\subsection*{Hour Breakdown}

\begin{center}
\begin{tabular}{lcc}
\toprule
\textbf{Task} & \textbf{Hours (est.)} & \textbf{Cost} \\
\midrule
Playing field definition & 3.5 & \$98.00 \\
Scoring system \& scorecards & 3.5 & \$98.00 \\
Student guide & 2.0 & \$56.00 \\
Moderator guide & 2.0 & \$56.00 \\
Anti-cheating documentation & 2.0 & \$56.00 \\
\midrule
\textbf{Phase 1 Total} & \textbf{13.0} & \textbf{\$364.00} \\
\bottomrule
\end{tabular}
\end{center}





\end{document}
